\scanpage{85}
% archive foliation: 84
% Printed face of the reused sheet; image rotated 180 degrees for reading.
% The upper limit of the initial sum lies beyond the torn edge and is not supplied.
\[
\sum_{j=0}(x_1T_1+\cdots+x_{k-1}T_{k-1})^j x_kT_k
(x_1T_1+\cdots+x_{k-1}T_{k-1})^{k-j-1}
\]
Etant donnés deux éléments $x,y$ de $L$, on définit $g_{i,k-i}(x,y)$ comme
le coefficient de $T_1^iT_2^{k-i}$ dans $(xT_1+yT_2)^k$.
Montrer que l'on a
\[
i!\ (k-i)!\ g_{i,k-i}(x,y)=f_k(t_1,\ldots,t_k)
\]
avec $t_h=x$ pour $h\leq i$ et $t_h=y$ pour $h\geq i+1$.

b) Lorsque $L$ est considérée comme algèbre de Lie sur $K$, on a, dans
$\mathcal L_K(L)$ :
\[
(\operatorname{ad}x)^n=\sum_{i=0}^n(-1)^{n-i}\binom ni L_x^i R_x^{n-i}.
\]
En déduire que, si $K$ est de caractéristique $p$ ($p$ premier), on a :
\[
\begin{aligned}
(1)\qquad &(\operatorname{ad}x)^p.y=(\operatorname{ad}x^p).y\\
(2)\qquad &(\operatorname{ad}x)^{p-1}.y=\sum_{i=0}^{p-1}x^iyx^{p-1-i}
\end{aligned}
\]
Déduire de (2) que :
\[
f_{p-1}(\operatorname{ad}x_1,\ldots,\operatorname{ad}x_{p-1}).y
=f_p(x_1,\ldots,x_{p-1},y)
\]
et :
\[
g_{i,p-1-i}(\operatorname{ad}x,\operatorname{ad}y).y
=(p-i)g_{i,p-i}(x,y).
\]
